% =====================================================================
%  CONCLUSIONES
% =====================================================================
\chapter*{Conclusiones}
\addcontentsline{toc}{chapter}{Conclusiones}

\begin{enumerate}
	\item El problema es real y cercano: en la UNSAAC, el 77,5\,\% de los estudiantes encuestados vivió al menos una manifestación de hostigamiento sexual y menos del 1\,\% lo denunció \citep{gutierrez2022}. Gran parte de lo que ocurre nunca llega a conocimiento de la universidad.
	\item Los antecedentes muestran que el aprendizaje automático ya se usa para detectar contenidos relacionados con el acoso, también en español. Sin embargo, esos trabajos usan publicaciones de redes sociales o mensajes; ninguno trabajó con relatos de estudiantes universitarios ni en el contexto peruano, y algunos se evaluaron de forma que puede sobrestimar sus resultados.
	\item CRISP-DM ofrece una ruta clara para construir el modelo, desde la recolección de relatos hasta el prototipo, y permite incorporar los cuidados éticos que exige trabajar con información sensible.
	\item El diseño experimental planteado permite comparar los modelos candidatos con las mismas condiciones y contrastar las hipótesis frente a una línea base léxica y una mayoritaria.
	\item El principal aporte para la Ingeniería Informática y de Sistemas será un modelo validado para detectar acoso sexual en relatos en español, integrado en un prototipo que apoya la revisión humana.
\end{enumerate}
